\documentclass[12pt,a4paper]{article}

% Essential packages
\usepackage[utf8]{inputenc}
\usepackage[T1]{fontenc}
\usepackage[margin=2.5cm]{geometry}
\usepackage{graphicx}
\usepackage{booktabs}
\usepackage{longtable}
\usepackage{array}
\usepackage{tabularx}
\usepackage[hidelinks]{hyperref}
\usepackage{natbib}
\usepackage{csquotes}
\usepackage{setspace}
\usepackage{titlesec}

% Font configuration
\usepackage{times}

% Line spacing
\onehalfspacing

% Section formatting
\titleformat{\section}{\Large\bfseries}{\thesection}{1em}{}
\titleformat{\subsection}{\large\bfseries}{\thesubsection}{1em}{}
\titleformat{\subsubsection}{\normalsize\bfseries}{\thesubsubsection}{1em}{}

% Bibliography style - numbered
\bibliographystyle{unsrt}
\setcitestyle{numbers,square}

% Custom commands for proper quotes
\newcommand{\openwindow}{``open window''}

\begin{document}

% Title page
\begin{titlepage}
\centering
\vspace*{2cm}

{\Huge\bfseries Scoping Review Protocol\par}
\vspace{1.5cm}

{\LARGE Mapping the Landscape of Precision Biomarkers for Exercise-Induced \openwindow: From Traditional Indices to Emerging Multi-Omics Candidates\par}
\vspace{2cm}

\hrule
\vspace{1cm}

{\large\bfseries Protocol Version:} 1.0\par
\vspace{0.5cm}
{\large\bfseries Date:} January 2026\par

\vspace{1cm}
\hrule

\vfill
\end{titlepage}

% Main content
\section{Background and Rationale}

\subsection{The ``Open Window'' Theory: Origins and Evolution}

Following high-intensity or prolonged exercise, transient alterations in immune function may occur, a phenomenon termed the ``exercise-induced open window'' \citep{nieman1994exercise}. This theory, first systematically articulated in the early 1990s, posits that athletes may experience a state of heightened immune susceptibility during the 3 to 72 hours following strenuous exercise, during which the risk of opportunistic infections, particularly upper respiratory tract infections (URTIs), is elevated. The immunological changes during this period are multifaceted, encompassing reductions in lymphocyte counts, suppression of natural killer (NK) cell function, alterations in neutrophil activity, and diminished mucosal immune function \citep{nieman1994exercise,gleeson2007immune}.

Early evidence supporting this theory derived primarily from epidemiological observations in endurance sports, marathon runners, long-distance runners, and cyclists demonstrated significantly higher infection rates following heavy training loads or competition compared to the general population \citep{peters1983ultramarathon,nieman1994exercise}. Notably, research examining the effects of resistance training, anaerobic exercise, or mixed exercise modalities on immune function remains relatively sparse, representing a gap in the current literature \citep{nieman2019compelling}. Existing evidence suggests that resistance exercise can also induce acute immune cell changes; however, whether these changes follow an \openwindow{} pattern similar to that observed in endurance exercise requires further investigation.

In recent years, the validity of this theory has been challenged by some investigators \citep{campbell2018debunking}. An alternative perspective proposes that the post-exercise decline in peripheral blood lymphocyte counts does not represent true immunosuppression, but rather reflects the redistribution of immune cells to peripheral tissues (such as mucosal surfaces, skin, and lungs). This redistribution may indicate enhanced immune surveillance rather than impaired immune function. Despite this ongoing debate, the dynamic changes in post-exercise immune status remain a central topic in exercise immunology, and accurate monitoring of immune status during this period holds significant practical implications for athlete health management.

\subsection{Traditional Immune Indices: Applications and Limitations}

Historically, exercise immunology research has relied upon a relatively fixed panel of traditional biomarkers to assess post-exercise immune alterations \citep{gleeson2007immune,nieman2019compelling}. These indices primarily include: peripheral blood leukocyte and lymphocyte subset counts, salivary secretory immunoglobulin A (sIgA) secretion rates, plasma cytokine concentrations (e.g., IL-6, IL-10, TNF-$\alpha$), serum cortisol levels, and natural killer (NK) cell counts and cytolytic activity. While these markers have established a foundational understanding of the relationship between exercise and immunity, their limitations have become increasingly apparent in the context of translating findings toward precision monitoring applications.

First, traditional indices represent single-dimensional measurements that inadequately capture the dynamic changes of the immune system as a complex network \citep{nieman2020future}. Immune responses involve the coordinated action of multiple cell types, signaling pathways, and effector molecules; single markers can only capture partial fragments of this complex landscape.

Second, traditional indices exhibit substantial inter-individual variability and lack clearly defined ``normal ranges'' or ``risk thresholds'' \citep{gleeson2007immune}. Baseline levels and exercise-response patterns for the same marker may differ considerably across individuals, rendering group mean-based interpretations inadequate for accurately reflecting the immune status of specific individuals. This limitation is particularly pronounced for precision monitoring approaches that aim to achieve individualized health management.

Third, the predictive association between traditional indices and actual infection risk has not been adequately validated \citep{campbell2018debunking,nieman2019compelling}. Although studies have observed changes in certain immune parameters following exercise, whether these changes truly translate into increased infection susceptibility currently lacks direct causal evidence. This uncertainty in the ``biomarker-outcome'' relationship constrains the clinical translational value of traditional indices.

Fourth, detection of traditional indices typically requires laboratory environments and specialized equipment, precluding convenient, real-time field monitoring, a limitation that conflicts with the demand for immediate feedback in athletic training settings.

These limitations collectively point to a core issue: the field of exercise immunology urgently requires exploration of biomarkers with greater precision, predictive capacity, and operational feasibility, namely, ``precision biomarkers'' capable of accurately reflecting individual immune status, predicting infection risk, and supporting individualized decision-making. This need has directed researchers' attention toward emerging multi-omics technologies.

\subsection{The Rise of Multi-Omics Technologies and the Current Fragmented Landscape}

In recent years, the development of high-throughput omics technologies has opened new perspectives for exercise immunology research \citep{nieman2020future}. Transcriptomics can reveal global changes in immune cell gene expression profiles following exercise; proteomics can identify protein networks associated with immune regulation in plasma or other biological fluids; metabolomics provides molecular-level information on immunometabolic interactions; and epigenetic studies (such as DNA methylation analysis) contribute to understanding the long-term effects of exercise on immune gene regulation. The application of these technologies has made the identification of novel biomarker candidates possible.

However, the current application of multi-omics in exercise immunology exhibits a distinctly fragmented character \citep{nieman2020future}. Existing studies predominantly focus on skeletal muscle adaptation or metabolic responses, rather than specifically targeting biomarker discovery for the post-exercise immune vulnerability period. The exercise protocols, sampling time points, and analytical platforms employed across different studies vary substantially, rendering results difficult to integrate. More importantly, these omics findings are primarily reported as ``research results'' rather than being systematically evaluated as potential precision monitoring tools.

\subsection{Research Gap and Methodological Rationale}

In summary, a clear gap exists in the current literature: no review has systematically integrated traditional immune indices with emerging multi-omics biomarker candidates within the specific framework of the ``exercise-induced open window'' to provide a comprehensive landscape mapping \citep{campbell2018debunking,nieman2019compelling,nieman2020future}. Existing reviews have either focused on theoretical debates regarding the hypothesis itself, addressed only single categories of biomarkers, or discussed omics technologies within the broader context of exercise physiology, none have achieved cross-level, cross-platform evidence integration.

Given the heterogeneity and dispersion of evidence in this field, the present study adopts a scoping review methodology. The core advantage of scoping reviews lies in their capacity to systematically map the evidence landscape of a research domain, identifying the distribution characteristics and gaps in existing research, without aiming for quantitative pooling of effect sizes \citep{arksey2005scoping}. This methodological characteristic aligns closely with the objective of this study to ``map the biomarker landscape.''

\section{Objectives and Research Questions}

\subsection{Primary Objective}

This scoping review aims to systematically map and synthesize the biomarkers documented in the existing literature for characterizing or monitoring immune status during the exercise-induced \openwindow{} period, encompassing the complete spectrum from traditional immunological indices to emerging multi-omics candidates, and to identify biomarkers or biomarker combinations with potential for precision monitoring applications.

It should be clarified that the objective of this study is descriptive rather than evaluative. This study does not seek to assess the diagnostic performance or predictive value of specific biomarkers, nor does it attempt to establish causal relationships between biomarkers and infection outcomes. The core task of this research is to: present the distribution characteristics of existing research in this field, reveal differences in research density across different biomarker types, and identify research directions that remain insufficiently explored.

\subsection{Research Questions}

To achieve the above objectives, this scoping review will address the following four specific research questions:

\textbf{Research Question 1:} What biomarkers have been used in the existing literature to characterize or monitor immune status changes during the post-exercise \openwindow{} period? Along what dimensions can these biomarkers be classified?

\textbf{Research Question 2:} To what extent have multi-omics technologies (transcriptomics, proteomics, metabolomics, epigenomics) been applied in this research context? What potential omics-derived candidate biomarkers have been identified?

\textbf{Research Question 3:} What correspondence exists between different biomarkers and sampling matrices (blood, saliva, urine, sweat, etc.)? What differences exist in the application frequency and research depth across different matrices?

\textbf{Research Question 4:} Based on existing evidence, which biomarkers or biomarker combinations have been considered by researchers to possess potential for translation into precision monitoring tools? How sufficient is the relevant evidence?

\section{Study Design}

This study employs a scoping review design, following the methodological framework proposed by Arksey and O'Malley \citep{arksey2005scoping} and its subsequent refinements \citep{levac2010scoping}, and is conducted in accordance with the Joanna Briggs Institute (JBI) guidelines for scoping reviews \citep{peters2020jbi}. The study report will be presented in accordance with the PRISMA-ScR (Preferred Reporting Items for Systematic Reviews and Meta-Analyses extension for Scoping Reviews) checklist \citep{tricco2018prisma}.

The selection of a scoping review rather than a systematic review or meta-analysis is based on the following methodological considerations:

First, the core objective of this study is to map the breadth and distribution of the evidence landscape, rather than to quantitatively pool effects for a specific intervention. The ``mapping'' function of scoping reviews directly corresponds to this objective.

Second, the literature anticipated for inclusion will encompass multiple study designs, multiple biomarker types, and diverse exercise protocols and outcome measures. This high degree of heterogeneity renders standardized effect size pooling neither feasible nor meaningful.

Third, this study aims to clarify how ``exercise-induced open window biomarkers'' have been conceptualized and operationalized in the existing literature, to identify research gaps, and to provide directional guidance for future research, all of which represent typical application scenarios for scoping reviews.

\section{Eligibility Criteria}

This study employs the PCC framework (Population-Concept-Context) recommended by JBI to construct inclusion and exclusion criteria, ensuring systematic and operational rigor \citep{peters2020jbi}.

\subsection{Population}

\textbf{Inclusion Criteria:}
This study includes literature with human subjects, specifically: healthy individuals participating in acute exercise tests or structured training programs; populations across different training levels, from untrained individuals and recreational exercisers to professionally trained athletes and elite competitors; no restrictions on sex or age.

\textbf{Exclusion Criteria:}
Studies primarily based on animal models or in vitro experiments are excluded. This exclusion is based on the following considerations: species differences exist between animal and human immune systems, limiting the extrapolation of findings to human athletes; this study focuses on human evidence with direct clinical translational potential. Studies with disease populations as primary subjects are also excluded, such as research using exercise as an intervention for cancer rehabilitation, cardiovascular disease management, or metabolic syndrome. Immune changes in such studies may be significantly influenced by underlying disease states, representing fundamental differences from the \openwindow{} phenomenon in healthy athletes.

\subsection{Concept}

\textbf{Inclusion Criteria:}
This study includes literature involving the following concepts:

\begin{itemize}
\item \textbf{Biomarkers:} Any measurable indicator used to reflect post-exercise immune status changes, including but not limited to cellular-level indices (leukocyte subset counts, NK cell activity), protein-level indices (cytokines, immunoglobulins), metabolic-level indices (metabolite profiles), and molecular-level indices (gene expression, epigenetic markers).

\item \textbf{Immune Monitoring:} Studies must involve the assessment or monitoring of exercise-related immune status, whether employing traditional methods or omics technologies.

\item \textbf{Exercise-Induced Open Window:} Studies must explicitly or implicitly address the post-exercise immune vulnerability period, immunosuppression period, or immune recovery period.
\end{itemize}

It should be specifically noted that this study does not require included literature to report actual infection outcomes (such as upper respiratory tract infection incidence). The research focus is on ``biomarkers associated with immune susceptibility'' themselves, rather than the predictive relationship between biomarkers and infection. This delineation ensures operational feasibility of the research scope while avoiding the omission of important biomarker discovery studies due to excessive restrictions.

\textbf{Exclusion Criteria:}
Studies focusing solely on exercise performance indices (such as lactate, maximal oxygen uptake) without involving immune-related outcomes are excluded; studies focusing on allergy or autoimmune diseases unrelated to infection susceptibility are excluded.

\subsection{Context}

\textbf{Inclusion Criteria:}
This study includes research conducted in the following contexts:

\begin{itemize}
\item Acute exercise protocols: Studies examining immune biomarker changes following a single exercise bout, with no restrictions on exercise modality (endurance, resistance, mixed, etc.).
\item Training studies: Studies monitoring dynamic changes in immune biomarkers during training cycles or training blocks.
\item Competition settings: Studies conducting immune monitoring in actual competition environments.
\end{itemize}

No restrictions are placed on the geographic location or institutional background of the research.

\textbf{Exclusion Criteria:}
Studies measuring only resting baseline immune parameters without involving exercise loads are excluded; studies with post-exercise sampling time points earlier than 30 minutes are excluded, as sampling time points that are too early cannot capture the dynamic immune changes during the \openwindow{} period.

\subsection{Types of Sources}

This scoping review will include the following types of original research: observational studies (cross-sectional, cohort, case-control), experimental studies (randomized and non-randomized controlled trials), and quasi-experimental designs. Published systematic reviews and narrative reviews will be used for reference tracking but will be recorded separately from original studies and not conflated. Conference abstracts, dissertations, and preprints will be excluded unless peer-reviewed full versions have been published.

\section{Search Strategy}

\subsection{Database Selection}

Systematic searches will be conducted in the following electronic databases: PubMed/MEDLINE (core biomedical database), Web of Science Core Collection (multidisciplinary coverage and citation tracking), Scopus (broad scientific literature coverage), SPORTDiscus (specialized sports science database), and Embase (supplementary biomedical coverage).

\subsection{Search Term Construction Logic}

The search strategy will be constructed around three core concept modules, with synonyms and related terms connected by ``OR'' within each module, and modules cross-limited by ``AND.''

\textbf{Concept Module One (Exercise):} Encompasses terms related to exercise, physical activity, training, athletic sports, endurance exercise, resistance training, high-intensity exercise, and related terminology.

\textbf{Concept Module Two (Immunity/Open Window):} Encompasses terms related to immunity, immune function, immunosuppression, open window, infection susceptibility, upper respiratory tract infection, mucosal immunity, and related terminology.

\textbf{Concept Module Three (Biomarkers/Omics):} Encompasses terms related to biomarkers, markers, cytokines, immunoglobulins, leukocytes, lymphocytes, NK cells, cortisol, transcriptomics, proteomics, metabolomics, epigenetics, multi-omics, and related terminology.

Different databases employ different search syntax (such as MeSH subject headings, Emtree terms, free-text words, etc.), but follow the same logical structure. Specific search strings will be optimized for each database with the assistance of large language models (LLMs).

\subsection{Supplementary Search}

In addition to database searches, this study will employ the following supplementary strategies: backward reference tracking of included literature; forward citation tracking of landmark publications in the field; and hand-searching of core journals such as \textit{Exercise Immunology Review} and \textit{Journal of Sport and Health Science}.

\subsection{Search Limitations}

Language is limited to English literature; the time range is set from January 2000 to the search execution date, to encompass research from the modern omics technology era while including foundational studies on traditional biomarkers.

\section{Study Selection Process}

\subsection{Screening Stages}

Study selection will be conducted in two stages:

\textbf{Stage One: Title and Abstract Screening.} All literature retrieved from database searches will undergo de-duplication, followed by independent preliminary screening at the title and abstract level by two reviewers (Reviewer A and Reviewer B). Screening will be based on the inclusion and exclusion criteria described above. Prior to formal screening, both reviewers will conduct a calibration exercise on a random sample of 50 articles to ensure consistent understanding of the inclusion criteria. Articles meeting inclusion criteria or those that cannot be determined based on title and abstract alone will proceed to full-text screening.

\textbf{Stage Two: Full-Text Screening.} Both reviewers will independently obtain and review the full texts of potentially eligible articles, making final inclusion decisions based on the complete inclusion and exclusion criteria. At this stage, specific reasons for excluding articles will be recorded and reported.

\subsection{Disagreement Resolution Mechanism}

During both screening stages, the screening results of Reviewer A and Reviewer B will be compared. For articles with discrepancies, both reviewers will first attempt to reach consensus through discussion. If agreement cannot be reached after discussion, a third reviewer will be introduced for arbitration. Inter-rater agreement during the screening process will be calculated and reported using Cohen's Kappa coefficient or percentage agreement.

\subsection{Process Documentation and Reporting}

The entire study selection process will be documented using reference management software (such as EndNote or Zotero) and screening management platforms (such as Rayyan or Covidence). A flow diagram will be generated in accordance with PRISMA-ScR requirements, clearly presenting the changes in article numbers at each stage and reasons for exclusion.

\section{Data Charting}

\subsection{Data Extraction Form}

A standardized data extraction form will be developed for this study to systematically record key information from included literature. The form will undergo pilot extraction on 10 articles prior to formal extraction, with iterative refinement based on pilot results.

Data extraction will encompass the following information domains:

\begin{longtable}{p{3.5cm}p{5cm}p{5.5cm}}
\toprule
\textbf{Information Domain} & \textbf{Extraction Items} & \textbf{Description} \\
\midrule
\endhead
\textbf{Study Characteristics} & First author and publication year & Study identification \\
 & Publication journal & Source tracking \\
 & Study design type & Observational/ Experimental/ Quasi-experimental \\
 & Country/Region of study & Geographic distribution \\
\midrule
\textbf{Participant Characteristics} & Sample size & Study scale \\
 & Demographic characteristics & Age, sex \\
 & Training status & Untrained/ Recreational/ Trained/ Elite athlete \\
 & Sport type & If applicable \\
\midrule
\textbf{Exercise Protocol} & Exercise type & Endurance/ Resistance/ Mixed/ Competition, etc. \\
 & Exercise intensity and duration & Load quantification \\
 & Study type & Acute exercise or training cycle study \\
\midrule
\textbf{Biomarker Information} & Biomarker name & Specific indicator \\
 & Biomarker category & Traditional immune index/ Cytokine/ Omics-derived, etc. \\
 & Omics level & Transcriptome/ Proteome/ Metabolome/ Epigenome \\
 & Sampling matrix & Blood/ Saliva/ Urine/ Other \\
 & Sampling time points & Time relative to exercise \\
 & Detection method & Analytical technique \\
\midrule
\textbf{Main Findings} & Direction and magnitude of biomarker change & Quantitative/ Qualitative description \\
 & Association with clinical outcomes & If reported \\
 & Author's assessment of biomarker utility & Translational potential evaluation \\
\bottomrule
\end{longtable}

\subsection{Extraction Process}

Data extraction will be completed independently by Reviewer A and Reviewer B. Both reviewers will extract data from the same batch of literature separately, followed by comparison. Items with discrepancies will be resolved through discussion, with verification against the original text when necessary. The specific process is illustrated in the figure below:

\begin{figure}[h]
\centering
\includegraphics[width=0.9\textwidth]{Figure.png}
\caption{Data Extraction Flowchart}
\label{fig:extraction}
\end{figure}

It should be emphasized that the purpose of data extraction is to serve the descriptive answering of research questions, rather than to provide a data foundation for subsequent statistical pooling. Therefore, the extraction focus is on biomarker types, distribution, and research context, rather than precise numerical effect sizes.

\section{Results Synthesis and Presentation Plan}

\subsection{Synthesis Method}

This scoping review will employ descriptive analytical methods to synthesize extracted data, rather than conducting quantitative meta-analysis. This choice aligns with the methodological positioning of scoping reviews: the objective is to present the distribution characteristics and research landscape of evidence, rather than to calculate pooled effect sizes.

Results synthesis will be organized around the four research questions, systematically addressing each through a combination of narrative description and visual presentation.

\subsection{Presentation Methods}

\textbf{Tabular Presentation:}

\begin{itemize}
\item Summary table of included studies: Presenting basic characteristics, study participants, exercise protocols, and biomarkers involved in each study
\item Biomarker inventory table: Systematically listing all identified biomarkers by category with their research frequencies
\item Sampling matrix-biomarker correspondence table: Displaying the correspondence between different biomarkers and sampling matrices
\end{itemize}

\textbf{Visual Presentation:}

\begin{itemize}
\item Biomarker-omics level mapping diagram: Presenting the complete spectrum from traditional indices to multi-omics candidates in conceptual map format, visually displaying research density at each level
\item Evidence temporal distribution diagram: Displaying the temporal evolution trends of research in this field, revealing the temporal transition points from traditional methods to omics approaches
\item Research gap matrix: Visually presenting biomarker-context combinations that have been adequately studied versus those that remain to be explored
\end{itemize}

\subsection{Synthesis Objectives}

The core objective of results synthesis is to reveal the distribution characteristics and gap areas of research in this field, rather than to compare or rank the merits of different biomarkers. The synthesis results will provide directional guidance for subsequent research and an evidence foundation of candidate biomarkers for the development of precision monitoring tools.

\section{Methodological Considerations and Quality Assessment Statement}

\subsection{Position on Formal Quality Assessment}

In accordance with the JBI scoping review methodological guidelines and PRISMA-ScR reporting standards \citep{peters2020jbi,tricco2018prisma}, scoping reviews typically do not require formal methodological quality assessment or risk of bias evaluation of included studies. This study will follow this methodological convention and will not conduct standardized quality scoring of included literature.

\subsection{Methodological Rationale}

This decision is based on the following considerations:

First, the core objective of scoping reviews is to map the scope and distribution of evidence, rather than to synthesize effects for a specific intervention. Under this objective, the manner in which study quality influences conclusions differs fundamentally from that in systematic reviews.

Second, the literature anticipated for inclusion will encompass highly heterogeneous study designs, making meaningful comparison using a unified quality assessment tool difficult.

Third, the JBI methodological manual explicitly states that scoping reviews ``do not aim to produce critically appraised and synthesized results/answers'' and that quality assessment ``is generally not performed'' \citep{peters2020jbi}.

\subsection{Handling of Quality-Related Information}

Although formal quality assessment will not be conducted, this study will still extract and report the following descriptive information related to study quality: study design type, sample size, whether validated detection methods were employed, appropriateness of sampling time points, and similar considerations. This information will be presented in the results section to serve as a reference for readers when interpreting the evidence.

% References
\bibliography{references}

\end{document}
